\documentclass[]{../../../../tex/classes/styledArticle}
\usepackage{../../../../tex/packages/hebrewSupport}


\author{שחר פרץ}
\title{\textit{היסטוריה 53}}
\begin{document}
	\maketitle
	
	בשאלה 8 במתכונת, יש לשים לב שהמקור מבטא את אחד הקשיים. ואז צריך להביא מקור מהספר. ומכאן, שיש צורך להבין אם יש מקור מתאים בספר לקושי שאתם רוצים לדבר עליו. 
	במקרים רבים כשיש כמה מטלות בסעיף הן קשורות אחת לשנייה ולכן חייבים לפני שניגשים לסעיף לדעת על מה כותבים, ובאיזה מקורות משתמשים. תבינו אם יש תלות בין החלקים ואז תתחילו לענות. 
	
	על הטופס למעלה כתוב כמה נספחים יש. תוודאו שיש לכם את כולם. 
	
	איפה לחפש חומרים בשירותים: 
	\begin{itemize}
		\item אם יש מראות, מאחורי המראות
		\item בתוך האסלה (תפתחו)
		\item בתוך ניילון במים
		\item מאחורי, מתחתי, מלמעלה
	\end{itemize}
	
	''תלמידים זה עם יצירתי``
	טיפקס לא, מרקר כהה לא, עט מחיק לא, דפים לא, שורה כן שורה לא. כתבו מספר שאלה ומספר סעיף. לא לכתוב בשני צידי הדף. בצד השמאלי את המחברת למעלה כתבו את תעודת הזהות. 
	
\end{document}